\documentclass[10pt]{article}
\usepackage[T1]{fontenc}
\usepackage[utf8]{inputenc}
\usepackage[french]{babel}
\usepackage{amsmath,amssymb,mathrsfs,mathpazo}
\usepackage{geometry}
\geometry{a4paper,margin=2cm}
\usepackage{graphicx,array,tabularx,tabularray,colortbl,multirow,multicol,booktabs,enumitem,pifont,ulem,fancybox,diagbox,cancel}
\usepackage{pstricks,pstricks-add,pst-plot,pst-tree,pst-node,pst-text,pst-func,pst-math,pst-bspline,pst-3dplot,pst-eucl,pst-solides3d}
\let\pscalculate\undefined
\usepackage{pst-bezier}
\usepackage{tikz,pgfplots,tkz-tab}
\usetikzlibrary{arrows,arrows.meta,patterns,calc,positioning,trees,decorations.markings,shapes.geometric,shapes.misc,shapes.arrows,intersections,decorations.pathmorphing,decorations.pathreplacing,backgrounds,shadows,babel}
\pgfplotsset{compat=1.18}
\usepgfplotslibrary{fillbetween}
\usepackage[np]{numprint}
\usepackage{listings,makecell,pas-tableur,twoopt}
\let\barre\undefined
\usepackage[pstricks]{tablvar}
\definecolor{darkgreen}{RGB}{0,100,0}
\providecommand{\red}{\color{red}}
\providecommand{\blue}{\color{blue}}
\providecommand{\green}{\color{green}}
\providecommand{\black}{\color{black}}
\providecommand{\white}{\color{white}}
\providecommand{\gray}{\color{gray}}
\providecommand{\cyan}{\color{cyan}}
\providecommand{\magenta}{\color{magenta}}
\providecommand{\yellow}{\color{yellow}}
\providecommand{\violet}{\color{violet}}
\providecommand{\orange}{\color{orange}}
\providecommand{\euro}{\texteuro}
\providecommand{\EUR}{\texteuro}
\providecommand{\ieme}{\textsuperscript{e}}
\providecommand{\degres}{\ensuremath{^\circ}}
\providecommand{\degree}{\ensuremath{^\circ}}
\providecommand{\vv}[1]{\overrightarrow{#1}}
\providecommand{\Coord}[1]{#1}
\providecommand{\decosix}{}
\providecommand{\starredbullet}{\ensuremath{\bullet}}
\pagestyle{empty}
\setlength{\parindent}{0pt}
\usepackage{adjustbox}\DeclareUnicodeCharacter{2236}{:}\DeclareUnicodeCharacter{2190}{\ensuremath{\leftarrow}}
\providecommand{\R}{\mathbb{R}}
\providecommand{\N}{\mathbb{N}}
\providecommand{\Z}{\mathbb{Z}}
\providecommand{\Q}{\mathbb{Q}}
\providecommand{\D}{\mathbb{D}}
\providecommand{\C}{\mathbb{C}}
\providecommand{\e}{{\mathrm{e}}}
\providecommand{\vect}[1]{\overrightarrow{#1}}
\providecommand{\vectt}[1]{\overrightarrow{\mathrm{#1}}}
\providecommand{\vvt}[1]{\overrightarrow{\mathrm{#1}}}
\providecommand{\np}[2][]{{#2~#1}}
\providecommand{\vv}[1]{{\overrightarrow{#1}}}
\renewcommand{\d}{{\mathrm{d}}}
\providecommand{\degres}{{^\circ}}
\providecommand{\degre}{{^\circ}}
\providecommand{\blue}{\color{blue}}
\providecommand{\red}{\color{red}}
\providecommand{\enskip}{\quad}
\providecommand{\llbracket}{\lbrack\!\lbrack}
\providecommand{\rrbracket}{\rbrack\!\rbrack}
\providecommand{\euro}{{\text{€}}}
\providecommand{\no}{{\text{n°}}}
\providecommand{\interval}[2]{{[#1;#2]}}
\providecommand{\intervFO}[2]{{[#1;#2[}}
\providecommand{\dotfill}{\ldots}

\providecommand{\ds}{\displaystyle}
\providecommand{\textding}[1]{#1}
\providecommand{\Oij}{\(\mathrm{O};\vec{\imath},\vec{\jmath}\)}
\providecommand{\Oijk}{\(\mathrm{O};\vec{\imath},\vec{\jmath},\vec{k}\)}
\providecommand{\Ouv}{\(\mathrm{O};\vec{u},\vec{v}\)}
\providecommand{\intervFF}[2]{\ensuremath{[#1;#2]}}
\renewcommand{\euro}{\ifmmode\text{\texteuro}\else\texteuro\fi}
\DeclareUnicodeCharacter{00B0}{\ifmmode^{\circ}\else\textdegree\fi}
\DeclareUnicodeCharacter{20AC}{\ifmmode\text{\texteuro}\else\texteuro\fi}
% Énoncés : Éducation nationale. Transcription et corrigés : APMEP, Denis Vergès et contributeurs. Découpage et classement : CoopMaths, Sébastien Lozano, Éric Elter et contributeurs. Adaptation de la présentation : Maths Explorer.
% CC BY-SA : https://coopmaths.fr/www/about/
% https://coopmaths.fr/alea/static/bac/2023/tex/bac_2023_03_sujet1_asie_3_cor.tex
\begin{document}
\noindent\textit{APMEP / CoopMaths. CC BY-SA. Présentation adaptée par Maths Explorer.}\par\medskip
\medskip


\begin{enumerate}
\item \textbf{Conjectures graphiques :}

Le nombre de solutions de l'équation $\ln (x) = kx$ :
\begin{itemize}
\item pour $k = 1$, l'équation ne semble pas avoir de solution, car il semble que la droite d'équation $y = x$ ne coupe pas la courbe représentative de la fonction ln;
\item pour $k = 0,2$, l'équation semble avoir deux solutions, car la droite d'équation $y=0,2x$ semble couper la courbe représentative de la fonction ln en deux endroits distincts.
\end{itemize}

\item Étude du cas $k = 1$ :
	\begin{enumerate}
		\item $f$ est dérivable sur $\R^{*+}$, en tant que différence de deux fonctions dérivables sur cet intervalle. Pour tout réel $x$ strictement positif, on a :

$f'(x) = \ln'(x) - 1 = \dfrac{1}{x} - 1 = \dfrac{1 - x}{x}$.

		\item On va donc étudier le signe de $f'(x)$ et en déduire les variations, au sein d'un tableau. Comme les limites aux bornes de l'intervalle de définition ne sont pas attendues, on ne va pas les déterminer.

Sur $\R^{*+}$, $x$ est strictement positif, donc le signe de $f'(x)$ est le signe de $1 - x$.

$1 - x > 0 \iff x < 1$. \quad $f(1) = \ln(1) - 1 = 0-1 = -1$.

On a donc :
\begin{center}






\psset{unit=1cm,arrowsize=2pt 3}
\begin{pspicture}(8.8,3.5)
\psframe(8.8,3.5)\psline(2.8,0)(2.8,3.5)\psline(0,2)(8.8,2)\psline(0,2.5)(8.8,2.5)
\psline(0,3)(8.8,3)\psline(2.9,0)(2.9,3)\psline(2.95,0)(2.95,3)
\uput[u](1.4,2.9){$x$} \uput[u](2.925,2.9){$0$} \uput[u](5.8,2.9){$1$} \uput[u](8.3,2.9){$+\infty$}
\rput(1.4,2.75){$1 - x$}\rput(1.4,2.25){$f'(x)$}\rput(1.4,1){variations de $f$}
\uput[u](4.4,2.4){+}\uput[u](5.8,2.4){0}\uput[u](7.3,2.4){$-$}
\uput[u](4.4,1.9){+}\uput[u](5.8,1.9){0}\uput[u](7.3,1.9){$-$}
\uput[d](5.8,2){$-1$}
\psline{->}(3,0.5)(5.3,1.5)\psline{->}(6.3,1.5)(8.3,0.5)
\end{pspicture}
\end{center}

\item On en déduit que sur son ensemble de définition, $f$ atteint un maximum égal à $-1$ pour $x = 1$, donc $f$ est à valeurs strictement négatives sur $\R^{*+}$.

Il n'y a donc aucune solution à l'équation $f(x) = 0$, or $f(x) = 0 \iff \ln(x) = x$, donc l'équation $\ln (x) = x$ n'admet aucune solution.
	\end{enumerate}
\item Étude du cas général :
	\begin{enumerate}
		\item Discutons du nombre de solutions :
\begin{itemize}
\item Si $g\left(\dfrac 1k\right)< 0$, alors la fonction $g$ a pour maximum un nombre réel strictement négatif, donc l'équation $g(x) = 0$ n'admet aucune solution.
\item Si $g\left(\dfrac 1k\right)= 0$, alors l'équation admet clairement $\dfrac{1}{k}$ comme solution et cette solution est la seule, car : $g$ étant strictement croissante sur $\left]0~;~\dfrac{1}{k}\right]$, $x < \dfrac{1}{k} \implies g(x) < g\left(\dfrac{1}{k}\right) $ et donc $g(x) < 0$.

Et $g$ étant strictement décroissante sur $\left[\dfrac{1}{k}~;~+\infty\right[$, $x > \dfrac{1}{k} \implies g(x) < g\left(\dfrac{1}{k}\right) $ et donc $g(x) < 0$.

Ainsi, l'équation admet une unique solution dans ce cas.

\item Si $g\left(\dfrac 1k\right)>0$, alors :

$f$ est continue (car dérivable) et strictement croissante sur $\left]0~;~\dfrac{1}{k}\right]$, et $0$ est une valeur intermédiaire entre $\lim\limits_{x \to 0}g = -\infty$ et $g\left(\dfrac{1}{k}\right)>0$, donc d'après le théorème de la bijection, l'équation $g(x) = 0$ admet une unique solution $\alpha$ sur cet intervalle.

Puis, de façon analogue, on applique le même théorème sur $\left]\dfrac{1}{k}~;~+\infty \right[$, où $g$ est continue et strictement décroissante pour faire émerger une seconde solution $\beta$, la seule sur $\left]\dfrac{1}{k}~~;~~+\infty \right[$.

Ainsi, l'équation $g(x) = 0$ admet exactement deux solutions dans ce cas.
\end{itemize}

		\item On rappelle que $k$ est un réel strictement positif, donc $\dfrac{1}{k}$ l'est aussi, et donc $\dfrac{1}{k}$ est dans l'ensemble de définition de $g$. On a $g\left(\dfrac 1k\right) = \ln\left(\dfrac 1k\right) - k\times \dfrac{1}{k} = -\ln(k) - 1$.

\item On a :
$\begin{aligned}[t] g\left(\dfrac 1k\right) > 0 &\iff -\ln(k) - 1 > 0\\
&\iff -\ln (k) >  1\\
&\iff \ln(k) < -1 \\ \end{aligned}$.

Nous parvenons bien à l'équivalence demandée.

\item On remarque que $g(x) = 0 \iff \ln(x) = kx$.

L'ensemble des valeurs de $k$ pour lesquelles l'équation $\ln (x) = kx$ possède exactement deux solutions est donc constitué des nombres $k$ vérifiant $\ln(k) < -1$, d'après la question précédente.

$\ln(k) < -1 \iff k < \e^{-1}$.

L'ensemble des valeurs $k$ cherché est donc $\left]0~;~\e^{-1}\right[ $
	\item En synthèse des questions précédentes, on peut donc dire que  l'équation $\ln (x) = kx$ :
			\begin{itemize}
				\item admet exactement deux solutions pour $ k\in \left]0~;~\e^{-1}\right[$;
				\item admet une unique solution pour $ k = \e^{-1}$;
				\item n'admet aucune solution pour $ k\in \left]\e^{-1}~;~+\infty\right[$;
			\end{itemize}
		\emph{Remarque :} Ce résultat confirme nos conjectures du début de problème, dans la mesure où $0,2 \in \left]0~;~\e^{-1}\right[$ et $1\in \left]\e^{-1}~;~+\infty\right[$ (en effet $\e^{-1}\approx 0,4)$.
		\end{enumerate}
	\end{enumerate}

	\medskip

\end{document}
